% chapter_05_Had.tex -- chapter 05 of the JitJet book: hadronisation (Lund strings), HadronLevel:Hadronize on, Decay off.
% Numbers: test/05_Had/out/ledger_05_Had.txt, genPartAnalyzer_05_Had.log, step1_GEN_05_Had.log, and the FWLite counts of
% doc/chapter/chapter_05_Had.md (section 4) read from step1_GEN_05_Had.root; the parton level of chapter 04 from
% test/04_MPI/out/ledger_04_MPI.txt.
\chapter{05 -- Hadronisation: the 169 partons become 320 primary hadrons}
\label{ch:05_Had}

\section{What this chapter is about}
Up to chapter~\ref{ch:04_MPI} the event is made of quarks and gluons. Nothing with a colour charge has ever
reached a detector: when the partons fly apart, the colour field between them does not spread out like an
electric field but is squeezed into a narrow tube, a \emph{string}, whose energy grows with its length until it
is cheaper to create a new quark--antiquark pair from the field than to stretch it further. This is
hadronisation, and in Pythia8 it is done by the Lund string model, switched on in this chapter with
\code{HadronLevel:Hadronize = on} while the decays of the hadrons stay off (\code{HadronLevel:Decay = off}).
The final state of this chapter is therefore the set of \emph{primary hadrons}: the particles the strings
produce directly, including the unstable ones ($\pi^0$, $\rho$, $\omega$, $K^*$, $\Delta$, $B^*$) that the next
chapter will decay. There are 320 of them, made from the 169 partons of chapter~\ref{ch:04_MPI}, and they carry
exactly the $13600.0\GeV$ of the two protons.\footnote{All numbers of this chapter are read from
\code{test/05\_Had/out/}: \code{ledger\_05\_Had.txt} (the record and the final state by kind and region),
\code{genPartAnalyzer\_05\_Had.log} (the $b$ and $\bar b$ chains, stage by stage), \code{step1\_GEN\_05\_Had.log}
(the Pythia8 settings and the hard process) and the FWLite counts of the companion
\code{doc/chapter/chapter\_05\_Had.md}, section~4, read from \code{step1\_GEN\_05\_Had.root}.}

The chapter follows the two partons of the hard process to the end of their colour strings: the $b$ quark
with $\pt=91.7\GeV$ becomes a $B^-$ meson with $\pt=94.3\GeV$, the $\bar b$ with $116.4\GeV$ a $B^{*+}$ with
$114.2\GeV$, both within $\Delta R<0.01$ of the quark. That a $b$ hadron keeps almost all of the momentum of
its quark is the single most important fact of this chapter for jet physics: it is what makes a $b$ jet
different from a light-quark or a gluon jet, and it is set by one fragmentation function and two numbers
that CMS has never tuned.

\section{From the previous chapter}
Chapter~\ref{ch:04_MPI} ended with 169 partons in the final state (status~1 in that file), counted there by the
scattering they belong to: the two $b$ quarks with the 12 ISR and FSR partons of the hard process, 142 partons of
the 31 additional multiparton interactions and their showers, and 13 beam-remnant quarks, all colour connected,
with $\sum E=13600.0\GeV$ and $\sum\pt=481.2\GeV$; the cone $\Delta R<0.4$ around the $b$ held $127.3\GeV$,
$6.9\GeV$ of it from the underlying event, the one around the $\bar b$ $117.1\GeV$. This chapter regenerates the same event with the same seeds,
so the parton level is identical: the same 169 partons are in this file with status~62 (156 outgoing
partons with their primordial $k_T$ included) and 63 (13 beam remnants), with the same $13600.0\GeV$ and
$481.2\GeV$, the same $b$ quark ($\pt=91.74\GeV$, $\eta=-1.57$) and $\bar b$ ($116.36\GeV$, $\eta=-3.75$).
Everything up to status~63 is kept. New in the record are 165 entries with status 71--79, the partons as they
enter the strings, reordered so that every colour singlet is contiguous (120 copies with status~71 and 45 partons
with status~73, each made by joining two very nearby ones), and the 320
primary hadrons: the record grows from 745 to 1230 entries. Gone is every free colour charge: the final
state of this chapter has 169 charged and 151 neutral hadrons and not a single parton.

\begin{plainwords}[Why hadronisation cannot be computed and must be modelled]
The parton shower stops at a scale of about $1\GeV$ (\code{TimeShower:pTmin}, raised to $0.611\GeV$ in the
log), where the strong coupling becomes too large for perturbation theory. From there to the hadrons there
is no formula, only a model with parameters fitted to data. Pythia8 uses the Lund string model, in which the colour field between a quark and an antiquark (or between the colour
ends of a chain of gluons) is a string with a constant tension $\kappa\approx1\GeV/\mathrm{fm}$. The string
breaks by creating $q\bar q$ pairs along its length, and each piece between two breaks becomes one hadron.
The model has a handful of parameters (which flavour is created at a break, how much momentum a hadron
takes, how much transverse momentum it gets); they were fitted once to $e^+e^-$ data at LEP, and CMS
inherits them unchanged through the tune CP5.
\end{plainwords}

\section{The numbers of the event}
\paragraph{The strings.}
Grouping the 320 primary hadrons by the set of partons they descend from gives 26 colour singlets, from
2 to 42 hadrons each (mean~12.3), listed in the companion. Twenty-one are ordinary strings with a quark at
one end and an antiquark at the other, from a plain $q\bar q$ pair up to chains of 10, 17 and 18 gluons
between the ends; two are \emph{junction} strings with three quark ends ($uud$ and $uds$), the leftovers of
the two proton remnants; three are \emph{ministrings}, $q\bar q$ pairs of mass $1.0$, $1.3$ and $1.5\GeV$, less
than \code{HadronLevel:mStringMin} $=1\GeV$ above the constituent masses of their two quarks, too light to
fragment and turned directly into two hadrons each. The most energetic string carries $3630.7\GeV$ along the
beam (a $u\bar d$ pair with 18 gluons, 42 hadrons); the heaviest has a mass of $1083.9\GeV$ (33 hadrons).
Table~\ref{tab:05strings} gives the two strings of the hard process: the $b$ quark is colour connected
through three gluons to an $\bar s$ quark created in the shower, a string of mass $12.8\GeV$ that gives six
hadrons; the $\bar b$ sits at the end of a plain two-parton string with a $d$ quark from the other side of
the event, of mass $99.6\GeV$, that gives only four hadrons because the $\bar b$ end keeps almost everything.

\begin{table}[htbp]
\centering\small
\caption{The two strings of the hard process and their primary hadrons (FWLite, \code{step1\_GEN\_05\_Had.root};
quark momenta of status~62, hadron momenta of the status-1 primary hadrons; $x_E=E(\text{hadron})/E(\text{quark})$).}
\label{tab:05strings}
\setlength{\tabcolsep}{4pt}
\begin{tabular}{lrrrrrl}
\toprule
string & quark \pt & quark $E$ & $M$ & $E$ & $n_{\rm had}$ & primary hadrons (GeV) \\
\midrule
$b\,g\,g\,g\,\bar s$ & 91.74 & 230.0 & 12.8 & 285.1 & 6 & $B^-$: $\pt=94.28$, $E=235.9$, $x_E=1.025$ \\
 & & & & & & $+\ \pi^0\,\pi^+\,\pi^0\,\pi^0\,K^0$ (\pt{} 0.9 to 7.2) \\
$d\,\bar b$ & 116.36 & 2469.9 & 99.6 & 2471.2 & 4 & $B^{*+}$: $\pt=114.17$, $E=2421.4$, $x_E=0.980$ \\
 & & & & & & $+\ \pi^-\,K^+\,K^{*-}$ (50 GeV between them) \\
\bottomrule
\end{tabular}
\end{table}

\paragraph{The fragmentation function and the $b$ hadrons.}
At each break the model draws the fraction $z$ of the remaining light-cone momentum of the string end that
the new hadron takes, from the Lund symmetric fragmentation function
\[
  f(z) \propto \frac{1}{z}\,(1-z)^{a}\,\exp\!\left(-\frac{b\,m_T^2}{z}\right),\qquad
  m_T^2 = m^2 + \pt^2 ,
\]
with $a=0.68$ (\code{StringZ:aLund}) and $b=0.98\GeV^{-2}$ (\code{StringZ:bLund}). For a light hadron
($m_T\approx0.5\GeV$) the exponential hardly matters and $f(z)$ peaks at small $z$: a light string end hands
out its momentum in many small steps, which is why the $u\bar d$ string of $3630.7\GeV$ gives 42 hadrons. For a
hadron containing the heavy quark at the string end, $m_T^2\approx m_B^2\approx28\GeV^2$ makes the exponential
kill all small $z$; in addition Pythia8 applies the Bowler modification for heavy quarks, an extra factor
$1/z^{\,r_Q\,b\,m_Q^2}$ with $r_b=0.855$ (\code{StringZ:rFactB}) and the $b$ mass $m_b=4.8\GeV$ of the hard
process. The result is a function that peaks at $z\approx0.93$ (mean $0.89$, for $m_T=m_B$) and vanishes at $z\to0$: the first
hadron produced at a $b$ end takes most of the string momentum, and the light hadrons that share the string
with it are soft. The event shows exactly this. The $B^-$ carries $x_E=1.025$ of the energy of its $b$ quark
and 83\% of the energy of the whole five-parton string ($235.9$ of $285.1\GeV$); it can carry \emph{more} than
its quark because the string also holds the three gluons and the $\bar s$ end, whose momentum is shared along
the string. The $B^{*+}$ carries $x_E=0.980$ of its $\bar b$ and leaves $50\GeV$ to three light hadrons. Both $b$
hadrons are within $\Delta R$ of $0.009$ and $0.001$ of their quark: hadronisation does not move the $b$
direction, it only dresses it. The five light hadrons of the $b$ string have $\pt$ between $0.9$ and
$7.2\GeV$.

\paragraph{The cones around the two $b$ hadrons.}
The dashed circles of figure~\ref{fig:05display} are the cones $\Delta R<0.4$ of chapters~\ref{ch:02_PS}
and~\ref{ch:04_MPI}, now around the $B^-$ and the $B^{*+}$, and the inset magnifies the cone of the $B^-$, whose
hadrons the main frame cannot separate. It holds nine primary hadrons with a vector \pt{} of $114.4\GeV$: the
$B^-$ and eight light hadrons with $\sum\pt=20.4\GeV$, five $\pi^0$ ($7.2$, $6.5$, $2.8$, $2.0$ and $0.5\GeV$), a
$K^0$, a $\bar K^{*0}$ and a $\pi^-$; the cone of the $B^{*+}$ holds a single $K^{*-}$ of $1.8\GeV$ besides it,
$116.0\GeV$. At parton level the same two cones held $127.3$ and $117.1\GeV$: hadronisation takes $12.8\GeV$
out of the $b$ cone. A cone is not a string: the hadrons are made along the strings, and several strings pass
through this cone (the $b$ string and those of the ISR and MPI partons that were inside it). Four hadrons just
outside the circle carry about as much as the cone lost: a $\bar p$ of $5.3\GeV$ at $\Delta R=0.45$, a
$\Delta^{++}$ of $2.8\GeV$ at $0.51$, and a $\rho^+$ and a $\pi^+$ of $2.4\GeV$ at $0.46$ and $0.42$, the $\pi^+$
one of the five light hadrons of the $b$ string.\footnote{\code{genPartAnalyzer\_05\_Had.log} (the ``cone'' line
of stage 4) and \code{ledger\_05\_Had.txt}; the hadrons outside the cone from FWLite (companion, section 4).}
This loss out of the cone is what the jet energy corrections later call the hadronisation part of the
out-of-cone showering; it is a property of the Lund strings, and of their parameters.

Which hadron is made at a break is the second draw of the model: a $q\bar q$ pair is created with
probability $u:d:s = 1:1:0.217$ (\code{StringFlav:probStoUD}), a diquark pair with probability 0.081 relative
to a quark pair (\code{StringFlav:probQQtoQ}, the source of the 34 baryons), and the new meson is a vector
rather than a pseudoscalar with a flavour-dependent ratio: $0.5$ for $u,d$ (\code{StringFlav:mesonUDvector}),
$0.55$ for $s$, $2.2$ for $b$ (\code{StringFlav:mesonBvector}). The last number is why the $\bar b$ became a
$B^{*+}$ and not a $B^+$: with a vector-to-pseudoscalar ratio of 2.2 the excited state is drawn 69\% of the
time, and the $B^{*+}$ will decay to $B^+\gamma$ in chapter~\ref{ch:06_Decay}.

\begin{plainwords}[Why a $b$ jet is not a light jet]
A light quark hands its momentum to a long chain of pions, each with a few per cent of it; a $b$ quark hands
80--100\% of it to one $B$ hadron in one step, because the Lund function with the Bowler term makes small $z$
impossible for a heavy end. A $b$ jet is therefore one heavy particle plus a few soft hadrons, and the jet
response of the detector, the $b$-tagging, the fraction of energy lost in neutrinos when the $B$ decays
(chapter~\ref{ch:06_Decay}) all follow from this one function and the two numbers $r_b$ and $b$, fitted to
$e^+e^-$ data at the $Z$ pole.
\end{plainwords}

\paragraph{The species.}
Table~\ref{tab:05species} counts the 320 primary hadrons: 286 mesons and 34 baryons, two of
them with a $b$ quark and two with a $c$ quark (a $\bar D^0$ and a $D^{*+}$ from a $c\bar c$ pair created in the
shower, with $\pt$ of only $2.0$ and $1.9\GeV$). The pions, charged and neutral, are 132; the vector mesons
$\rho$ and $\omega$, 77, are about half of the 149 light pseudoscalars ($\pi$, $\eta$, $\eta'$), which is the
vector-to-pseudoscalar ratio of 0.5 at work (one spin state for the pseudoscalar against three for the vector,
suppressed by their mass); the
kaons and $K^*$, 56, reflect the strangeness suppression 0.217. The spectrum is soft: 124 hadrons have
$\pt<0.5\GeV$, 226 have $\pt<1\GeV$, only 8 have $\pt>5\GeV$ and only the two $B$ hadrons exceed $10\GeV$.
The transverse momentum of a light hadron comes from the Gaussian $\pt$ kick of each string break, of width
$\sigma=0.335\GeV$ (\code{StringPT:sigma}); this is where the $+24.5\GeV$ of scalar $\sum\pt$ of the ledger
comes from, $0.16\GeV$ per additional particle.

\begin{table}[htbp]
\centering\small
\caption{The 320 primary hadrons by class (FWLite count of the status-1 particles of
\code{step1\_GEN\_05\_Had.root}, listed by species in the companion, section~4).}
\label{tab:05species}
\begin{tabular}{lrrrl}
\toprule
class & $n$ & $\sum E$ [GeV] & $\sum\pt$ [GeV] & species \\
\midrule
pions & 132 & 2736.5 & 108.5 & $\pi^\pm$ 86, $\pi^0$ 46 \\
light vector and $\eta$ mesons & 94 & 3169.2 & 99.0 & $\rho^\pm$ 34, $\rho^0$ 25, $\omega$ 18, $\eta$ 14, $\eta'$ 3 \\
kaons & 39 & 1173.6 & 39.0 & $K^\pm$ 22, $K^0$ 17 \\
$K^*$ & 17 & 414.8 & 13.1 & $K^{*0}$ 9, $K^{*\pm}$ 8 \\
nucleons & 19 & 2888.6 & 18.0 & $n$ 10, $p$ 9 \\
other light baryons & 15 & 523.8 & 15.7 & $\Lambda$ 3, $\Sigma$ 3, $\Sigma^*$ 4, $\Delta$ 5 \\
charm hadrons & 2 & 36.2 & 3.9 & $\bar D^0$, $D^{*+}$ \\
bottom hadrons & 2 & 2657.3 & 208.5 & $B^-$, $B^{*+}$ \\
\midrule
all & 320 & 13600.0 & 505.7 & 169 charged, 151 neutral \\
\bottomrule
\end{tabular}
\end{table}


\paragraph{Energy is conserved exactly.}
The ledger of chapter~\ref{ch:04_MPI} closed at $13600.0\GeV$ and this one closes at $13600.0\GeV$: not
approximately, but string by string. String fragmentation takes the total four-momentum of each colour
singlet and distributes it among its hadrons; the last two hadrons of every string are made together from
whatever four-momentum is left (the ``final joining''), so that each string conserves energy and momentum
by construction. In this event the 26 strings have $|\Delta E|<0.5\MeV$ each between their partons and
their hadrons. What hadronisation does create is mass: the 169 partons had $\sum m=30.8\GeV$ ($9.6\GeV$ in the
two $b$ quarks, $16.7\GeV$ in the 39 other quarks of the showers and the MPI, $4.5\GeV$ in the 13 remnant
quarks; the gluons are massless), the 320 hadrons have $\sum m=176.7\GeV$, taken from the
kinetic energy of the strings. Pythia8 then checks the sum over the whole event and would reject it if the
mismatch exceeded \code{Check:epTolErr}, set to $0.01$ by CMS in
\path{Configuration/Generator/python/Pythia8CommonSettings_cfi.py} (Pythia8 default $10^{-4}$); the check never
fires here.

\paragraph{Colour reconnection.}
Before the strings are formed, Pythia8 may rearrange the colour connections between the 32 scattering systems
of the event (the hard process and 31 MPI): this is colour reconnection (CR), and in the MPI-based model used by
CP5 (\code{ColourReconnection:mode = 0}) a system of scale $\pt$ is merged into a harder one with probability
$(p_{\mathrm{T},0}^{\mathrm{rec}})^2/((p_{\mathrm{T},0}^{\mathrm{rec}})^2+\pt^2)$, with $p_{\mathrm{T},0}^{\mathrm{rec}}=\mathrm{range}\times p_{\mathrm{T},0}$. CP5 sets
\code{ColourReconnection:range = 5.176} (Pythia8 default 1.8, the largest change CP5 makes to the hadron level);
with the $p_{\mathrm{T},0}=1.44\GeV$ of this run (\path{step1_GEN_05_Had.log}) a system of $\pt\approx7.5\GeV$ is merged
half of the time, and the soft MPI systems almost always. CR shortens the total string length and so lowers
the number of hadrons per MPI; it is the knob with which the tune reproduces the measured $\langle\pt\rangle$
versus charged multiplicity. It is why the 32 systems of this event end up as 26 strings, several with ten
and more gluons from different scatterings. It leaves no separate status code in the record (status~79
belongs to the other CR models): its effect is in the colour indices of the status-62 partons.

\section{Ledger and display}
Table~\ref{tab:05ledger} is the ledger of the stage: the 169 partons before (identical to the final state of
chapter~\ref{ch:04_MPI}, and present in this file as status~62 and 63) against the 320 primary hadrons after.
The energy is unchanged in total but moves between regions: the forward strings span many units of
rapidity and their hadrons land where the string, not the parton, was. The region beyond $\abseta=5.5$ loses
$604.6\GeV$ and the ring $5.2<\abseta<5.5$ gains $565.4\GeV$ (fifteen hadrons there instead of three
partons), while the tracker region gains 61 particles and $42.5\GeV$. The visible energy within $\abseta<5.2$
goes from $5740.0$ to $5779.2\GeV$, less than 1\% of a change; the scalar $\sum\pt$ rises by 5\% for the
reason given above. The two cone rows are the jet side of the ledger: hadronisation takes $12.8\GeV$ of \pt{}
out of the $b$ cone and $1.1\GeV$ out of the $\bar b$ cone. Figure~\ref{fig:05display} is the $\eta$--$\phi$
display of the primary hadrons: the $B^-$ at $(-1.57,-2.09)$ and the $B^{*+}$ at $(-3.75,1.01)$ are the only two
large markers, each in its cone; the rest is the soft spray of the 26 strings, denser at large $\abseta$, with
25 hadrons beyond $\abseta=6$ drawn on the border.

\begin{plainwords}[What a detector will see of this stage]
Nothing of the strings themselves, only the hadrons, and only after the next chapter has decayed the unstable
ones. But the split made here decides what kind of energy arrives: 169 charged hadrons with $355\GeV$ of \pt{}
for the tracker, 151 neutral ones with $151\GeV$, among them 46 $\pi^0$ that will become photons in the
electromagnetic calorimeter. The jet energy scale is the detector response to this mixture, and the mixture is
set by the Lund parameters.
\end{plainwords}

\begin{center}\begin{minipage}{\textwidth}\centering\footnotesize
\captionof{table}{Ledger of the hadronisation stage, event 1:1:1. Before: the 169 partons of chapter~\ref{ch:04_MPI}
(\code{test/04\_MPI/out/ledger\_04\_MPI.txt}, final state), after: the 320 primary hadrons
(\code{test/05\_Had/out/ledger\_05\_Had.txt}); the $b$ and $\bar b$ rows from \code{genPartAnalyzer\_05\_Had.log},
the masses from the FWLite count of the companion. $E$ and \pt{} in GeV, \pt{} scalar sums, except in the cone
rows: the final-state particles within $\Delta R<0.4$ of the last copy of the $b$ and $\bar b$ quark
(\code{genLedger.py}), vector \pt; differences before rounding.}
\label{tab:05ledger}
\begin{tabular}{lrrrrrrrrr}
\toprule
 & \multicolumn{3}{c}{before (partons)} & \multicolumn{3}{c}{after (primary hadrons)} & \multicolumn{3}{c}{difference} \\
\cmidrule(lr){2-4}\cmidrule(lr){5-7}\cmidrule(lr){8-10}
 & $n$ & $\sum E$ & $\sum\pt$ & $n$ & $\sum E$ & $\sum\pt$ & $n$ & $\sum E$ & $\sum\pt$ \\
\midrule
final state & 169 & 13600.0 & 481.2 & 320 & 13600.0 & 505.7 & $+151$ & $0.0$ & $+24.5$ \\
\quad partons & 169 & 13600.0 & 481.2 & 0 & 0.0 & 0.0 & $-169$ & $-13600.0$ & $-481.2$ \\
\quad charged hadrons & 0 & 0.0 & 0.0 & 169 & 8863.8 & 355.0 & $+169$ & $+8863.8$ & $+355.0$ \\
\quad neutral hadrons & 0 & 0.0 & 0.0 & 151 & 4736.2 & 150.7 & $+151$ & $+4736.2$ & $+150.7$ \\
\midrule
$\abseta<2.5$ (tracker) & 80 & 642.1 & 254.5 & 141 & 684.6 & 259.9 & $+61$ & $+42.5$ & $+5.4$ \\
$2.5<\abseta<3.0$ & 17 & 190.5 & 22.3 & 30 & 196.8 & 24.6 & $+13$ & $+6.3$ & $+2.3$ \\
$3.0<\abseta<5.2$ (HF) & 50 & 4907.4 & 183.0 & 94 & 4897.9 & 192.4 & $+44$ & $-9.5$ & $+9.4$ \\
$5.2<\abseta<5.5$ & 3 & 245.3 & 2.5 & 15 & 810.7 & 7.9 & $+12$ & $+565.4$ & $+5.4$ \\
$\abseta>5.5$ & 19 & 7614.7 & 18.9 & 40 & 7010.1 & 20.9 & $+21$ & $-604.6$ & $+2.0$ \\
visible, $\abseta<5.2$ & 147 & 5740.0 & 459.8 & 265 & 5779.2 & 476.9 & $+118$ & $+39.2$ & $+17.1$ \\
\midrule
$\sum m$ of the final state & & 30.8 & & & 176.7 & & & $+145.9$ & \\
$b$ quark $\to$ $B^-$, \pt & & & 91.74 & & & 94.28 & & & $+2.54$ \\
$\bar b$ quark $\to$ $B^{*+}$, \pt & & & 116.36 & & & 114.17 & & & $-2.18$ \\
$|\vec\pt(b)+\vec\pt(\bar b)|$ & & & 25.08 & & & 20.30 & & & $-4.77$ \\
\midrule
$b$ cone $\Delta R<0.4$ & 9 & 325.3 & 127.3 & 9 & 290.6 & 114.4 & 0 & $-34.8$ & $-12.8$ \\
$\bar b$ cone $\Delta R<0.4$ & 2 & 2485.9 & 117.1 & 2 & 2465.3 & 116.0 & 0 & $-20.6$ & $-1.1$ \\
\bottomrule
\end{tabular}
\end{minipage}\end{center}

\begin{figure}[htbp]
\centering
\includegraphics[width=\textwidth]{05_Had/out/genPartAnalyzer_05_Had_page-4.png}
\caption{The event after hadronisation (\code{test/05\_Had/out/genPartAnalyzer\_05\_Had\_page-4.png}, page 4 of
\code{genPartAnalyzer\_05\_Had.pdf}): the 320 primary hadrons in $\eta$--$\phi$, marker size $\propto\log\pt$,
the $B^-$ (red) and the $B^{*+}$ (blue) matched to the $b$ and $\bar b$ chains, the star their vector sum. Dashed:
the cones $\Delta R<0.4$ (analyzer option \code{coneR}), $114.4\GeV$ around the $B^-$ (nine hadrons) and
$116.0\GeV$ around the $B^{*+}$ (two); the inset magnifies the $B^-$ cone with the name and \pt{} of every hadron
within $1.25\times0.4$ of its axis (``+4 more'': the softest, not labelled). Page 3 of the same file is the end of
the parton level, the display of chapter~\ref{ch:04_MPI}, with the quarks at the same places.}
\label{fig:05display}
\end{figure}


\FloatBarrier
\section{Open problems for Phase-2}
\paragraph{The hadronisation of CP5 is Pythia8 as delivered.}
The tune of this event, CP5, is the file
\path{Configuration/Generator/python/MCTunes2017/PythiaCP5Settings_cfi.py} of the release, 21 settings: \code{Tune:pp 14}
and \code{Tune:ee 7} (the Monash 2013 tune as base), the MPI parameters (\code{pT0Ref=1.41}, \code{ecmPow=0.03344},
\code{coreRadius}, \code{coreFraction}), \code{ColourReconnection:range=5.176}, $\alpha_S=0.118$ at second order for
the showers and the MPI, the PDF \code{NNPDF31\_nnlo\_as\_0118} and the total cross sections. It does not contain a
single \code{StringZ}, \code{StringPT} or \code{StringFlav} parameter, nor does \path{Pythia8CommonSettings_cfi.py}
(\code{Check:epTolErr = 0.01} and the decay settings of chapter~\ref{ch:06_Decay}). The fragmentation of every CMS
Pythia8 sample is therefore the Pythia8 default, which in the Pythia 8.311 of this release
(\path{$PYTHIA8DATA/Fragmentation.xml}, \path{FlavourSelection.xml}) is \code{StringZ:aLund = 0.68}, \code{bLund = 0.98},
\code{rFactB = 0.855}, \code{rFactC = 1.32}, \code{aExtraDiquark = 0.97}, \code{StringPT:sigma = 0.335},
\code{StringFlav:probStoUD = 0.217}, \code{probQQtoQ = 0.081}, \code{mesonBvector = 2.20}; the ``changes only'' table
that Pythia8 prints in \path{step1_GEN_05_Had.log} lists none of them. These are the Monash 2013 values fitted to LEP
and SLD $e^+e^-$ data. \code{rFactB}, the Bowler parameter that fixes how much of the $b$ momentum the $B$ hadron
takes, was fitted to the $B$-hadron energy fraction measured at the $Z$ pole and has never been set by CMS: not in
CP5, not in its \code{TuneUp}/\code{TuneDown} variants (which vary only the MPI and CR parameters), not in the Run-3
copy \path{MCTunesRun3ECM13p6TeV/PythiaCP5Settings_cfi.py} (which differs from the 2017 file only in
\code{SigmaTotal:sigmaEl = 22.08}, \code{sigmaTot = 101.037}). The only tune files that touch any string parameter
are \path{MCTunes2017/PythiaCP5CR1TuneSettings_cfi.py}, the colour-reconnection variant for the top-mass systematics
(\code{ColourReconnection:mode=1} with \code{StringZ:aLund=0.38}, \code{bLund=0.64}, \code{probQQtoQ=0.078},
\code{probStoUD=0.2}), and its $13.6\TeV$ copy; their sibling \code{CP5CR2} (\code{mode=2}) changes the colour
reconnection only. Outside the tune files a single fragment of \path{Configuration/Generator/python} sets a string
parameter: the Run-1 \path{UpsMM_8TeV_TuneCUETP8M1_cfi.py}, with \code{StringFlav:mesonBvector = 3.000} and the
comment ``needs work''. Since when: the directory is \code{MCTunes2017} and CP5 dates
from the 2017 tuning campaign (CMS-GEN-17-001); the release on \code{cvmfs} carries no git history, so the date at
which the file took its present content could not be established here.

\paragraph{Hadronisation does not know the collider energy, the jets do.}
None of the string parameters depends on $\sqrt s$: a $b$ string of $12.8\GeV$ fragments the same way at
13.6 and at $14\TeV$, and the Phase-2 relvals use the same tune file: \path{TTbar_14TeV_TuneCP5_cfi.py},
\path{QCD_Pt15To7000_Flat_14TeV_TuneCP5_cff.py} and \path{MinBias_14TeV_pythia8_TuneCP5_cfi.py} (the fragments of
\path{Configuration/PyReleaseValidation/python/upgradeWorkflowComponents.py}) all import
\path{MCTunes2017.PythiaCP5Settings_cfi}; there is no \code{MCTunes} directory for $14\TeV$. What changes at the
HL-LHC is everything around this stage: at $14\TeV$ the flavour composition of the jets shifts towards gluon
jets, which fragment into more and softer hadrons (the 20-parton string of this event gives 42 hadrons, the
$b$ string 6), and the response of a particle-flow detector depends on that multiplicity and on the
charged-to-neutral split (169 to 151 here). The colour reconnection of CP5 is tuned to the Run-2 underlying
event, and \code{ColourReconnection:range} enters through $p_{\mathrm{T},0}$, which CP5 scales as
$(\sqrt s/7\TeV)^{0.03344}$ (\code{MultipartonInteractions:ecmPow}); with 200 pileup collisions per crossing,
each hadronised separately with the same strings, an error in the soft hadron spectrum ($\sigma$, $a$, $b$)
under every jet is multiplied by 200.

\paragraph{What the Phase-2 jet-energy calibration will need from this stage.}
Three things, none of which is in a configuration file today. (i) Flavour-dependent corrections and their
uncertainty: $b$, light-quark and gluon jets respond differently because of the fragmentation function; the
standard estimate of the uncertainty is the difference to an independent hadronisation model, which the release
has (\path{Configuration/Generator/python/Herwig7Settings/Herwig7CH3TuneSettings_cfi.py}, the cluster model of
Herwig7), but no Phase-2 relval workflow with Herwig7 exists in \path{upgradeWorkflowComponents.py}. (ii) A $b$-jet energy scale with a $b$-fragmentation uncertainty: the
fraction $x_E$ of the $b$ momentum in the $B$ hadron is governed by \code{rFactB = 0.855}, a LEP number, whose
variations ($r_b$, $b$) have to be propagated to the $b$-jet response together with the semileptonic branching
fractions of chapter~\ref{ch:06_Decay}; with the Phase-2 tracker to $\abseta=4$ and the MTD timing layer, $b$
jets will be tagged over a wider region than today. (iii) The forward region: $94$ primary hadrons with
$4897.9\GeV$ sit in $3.0<\abseta<5.2$, where Phase-2 puts HGCal and a tracker up to $\abseta=4$ instead of the
HF alone; the multiplicity and neutral fraction the strings give there have never been measured by CMS and
are an extrapolation of the central tune. Whether the jet-energy-correction payloads of the Phase-2 global
tags carry flavour-dependent corrections lives in the conditions database, not in the python of the release,
and was not checked here.
